% Section fragment for the parent article; requires amsmath, amssymb,
% amsthm, and theorem/lemma/proposition/corollary environments.
\section{Few nonlinear columns: an exact fixed-parameter forest algorithm}
\label{sec:fpt-column-theorem}

Let $R$ be the row set of the immutable nonnegative occurrence matrix,
and let $E\subseteq R$ label all nonloop edges of a connected undirected
multigraph $G$. Parallel edges are allowed; other rows are matroid loops. A selected forest
activates its rows initially. A unit-column pivot activates its last
remaining row. We use the established equivalence between complete
positive raw cancellation and complete unit-column peeling for this
fixed source presentation.

An active column has at least one entry equal to one; a column without
such an entry is inert for unit peeling. Form the ordinary implication
digraph $P$ on $R$ from all active columns supported on at most two rows.
A singleton unit column contributes its row to the seed set $Z$.
A two-row column supported on $u,v$ contributes $v\to u$ if its entry
at $u$ is one, and $u\to v$ if its entry at $v$ is one. Thus a
$(1,1)$ column contributes both directions and a $(1,a)$ column with
$a\ge2$ contributes only the direction towards the unit entry.

Let $H_1,\ldots,H_k$ be the supports of the remaining active columns.
Assume each consists of three distinct unit owners. Each such column
contributes the three implications
\[
                  H_i\setminus\{x\}\longrightarrow x
                  \qquad(x\in H_i).
\]
Ordinary implication closure is directed reachability in $P$, denoted
$\operatorname{Reach}_P$. The full activation closure additionally uses
these two-premise implications. An exceptional column can produce at
most one newly activated row: after its first useful application, all
three of its support rows are active.

\subsection{Removing the closure that is free for every tree}

Write $\cl_A(X)$ for the rows activated by unit peeling
from an initial row set $X$, including all singleton-unit seeds. Compute
the \emph{free closure}
\[
                         W=\cl_A(\varnothing).
\]
This preprocessing involves no choice of top forest. Restrict the
matrix columns to the live rows $R_0=R\setminus W$, keeping their
coefficients unchanged.

\begin{proposition}[Exact free-closure preprocessing]
\label{prop:fpt-free-closure}
For every initially selected row set $X\subseteq R$,
\[
 \cl_A(X)
   =W\cup
     \cl_{A[R_0]}(X\setminus W).
\]
Consequently the exceptional-column parameter may be measured after
this preprocessing. The original cap graph and every original eligible
edge must remain available for tree completion and weight optimization.
\end{proposition}
\begin{proof}
The activation rule for a unit entry in row $x$ requires all other
support rows of its column to have been activated. Every row of $W$ is
active from every initial set $X$, by monotonicity. Once these rows are
active, the rules with outputs in $R_0$ have precisely the prerequisites
obtained by restricting their columns to $R_0$. Each residual activation
is therefore valid in the original system after activating $W$.
Conversely, remove from an original activation schedule the events whose
outputs lie in $W$. Every remaining event has all of its remaining
prerequisites activated earlier in the residual schedule, giving the
reverse inclusion. A column used during free peeling has no live owner
afterwards, so it cannot create a conflict through reuse.
\end{proof}

One convenient implementation retains all original row labels, marks
the rows of $W$ as free seeds, and removes them only from the residual
column supports. Their vertices then have no residual clauses. This
keeps the original connected graph and its eligible edge set unchanged.
Equivalently, discarded free rows may be retained as extension-only graph
edges. A column losing one of its three unit owners becomes an ordinary
two-row column, so the residual value of $k$ cannot increase. Authenticate
any geometric parity identity from the uncapped source counts before
capping multiplicities.

Henceforth, recompute $P,H_1,\ldots,H_k$ from the residual columns,
retain the original $R,E,G$, and include every row of $W$ in $Z$.
The rows of $W$ are isolated free seeds in the residual implication
system. All subsequent parameters refer to this recomputed system.

\begin{theorem}[Fixed-parameter forest selection by exceptional columns]
\label{thm:column-fpt}
The existence of a spanning tree $T\subseteq E$ whose complementary
rows peel completely can be decided deterministically in
\[
       3^k k! B_{2k}\,\operatorname{poly}(N)
       =2^{O(k\log(k+2))}\operatorname{poly}(N)
\]
bit operations, where $N$ is the encoded graph--matrix input length and
$B_j$ is the $j$th Bell number. A successful tree and an immutable-source
peeling schedule can be constructed within this bound. The algorithm
enumerates no $|R|^k$ family of forest-edge choices.
\end{theorem}

The parameter in this theorem is the number of exceptional
\emph{columns}, and the algorithm chooses the forest jointly with the
peeling strategy. It is therefore distinct from an algorithm that fixes
the forest first or parameterizes exceptional rows. Its proof follows
from a reduction to classical independent transversals; no matroid
property is assumed for the family of peelable row sets.

\subsection{Ordered exceptional events}

A \emph{scenario} consists of a permutation $i_1,\ldots,i_k$ of the
exceptional columns and an output choice $o_j\in H_{i_j}$ for each
position. At event $j$, the two rows of $H_{i_j}\setminus\{o_j\}$
must already be active, and $o_j$ becomes active. An output already
active is allowed; the event then has no effect. There are $3^k k!$
scenarios, including one empty scenario when $k=0$.

\begin{lemma}[Complete scenario family]
\label{lem:fpt-scenarios-complete}
Every successful forest admits a complete activation schedule obeying
one of these scenarios. Conversely, a legal scenario execution that
activates every row certifies success of its initial forest.
\end{lemma}
\begin{proof}
Take a complete activation schedule and list, in their actual order,
the exceptional-column applications that activate a new row. No column
appears twice in this list. Append every unused exceptional column after
the original schedule finishes, choosing its output arbitrarily. These
appended events are legal because all rows are already active. The
converse follows directly from the implication rules: each useful
exceptional event is a column with exactly one unit occurrence remaining.
Ordinary directed reachability is implemented by the corresponding
two-row or singleton unit columns.
\end{proof}

Fix a scenario and write $O=\{o_1,\ldots,o_k\}$. Let
$Q_1,\ldots,Q_p$ be those source strongly connected components of $P$
that meet neither $Z$ nor $O$, and put
\[
              A_\ell=Q_\ell\cap E,\qquad
              U=Z\cup\bigcup_{\ell=1}^p Q_\ell.
\]
The sets $A_\ell$ are pairwise disjoint. Every complete execution of
this scenario must select a row in each $A_\ell$: no ordinary arc enters
$Q_\ell$, and no chosen exceptional event outputs into it. Conversely,
one selected row in $Q_\ell$ activates all of $Q_\ell$ by strong
connectivity. These are obligations \emph{conditioned on the scenario};
they need not be stopping sets of the full matrix when exceptional
columns are present.

For a prerequisite $x\in H_{i_j}\setminus\{o_j\}$ at event $j$, set
\[
 C_j=\operatorname{Reach}_P
         \bigl(U\cup\{o_1,\ldots,o_{j-1}\}\bigr).
\]
If $x\in C_j$, no additional requirement is needed. Otherwise introduce
the \emph{predecessor demand}
\[
 D_{j,x}=\{e\in E:
              x\in\operatorname{Reach}_P(\{e\})\}.
 \tag{$\ast$}\label{eq:fpt-predecessor-demand}
\]
There are at most $2k$ retained demands. Repeated demands may be kept as
separately indexed requirements. Only earlier outputs occur in $C_j$;
using future outputs would permit circular activation justifications.

\begin{lemma}[Separation of the residual demands]
\label{lem:fpt-demand-separation}
Every retained predecessor demand is disjoint from every $A_\ell$.
\end{lemma}
\begin{proof}
If $e\in A_\ell\cap D_{j,x}$, then $e\in Q_\ell\subseteq U$ and
there is a directed path from $e$ to $x$. Consequently $x\in C_j$, in
contradiction to retention of this demand.
\end{proof}

\begin{proposition}[Exact requirements for a fixed scenario]
\label{prop:fpt-fixed-scenario}
An independent edge set $F$ supports complete activation according to the
fixed scenario if and only if $F$ meets every $A_\ell$ and every retained
predecessor demand.
\end{proposition}
\begin{proof}
For necessity, the source-component argument gives the requirements
$F\cap A_\ell\ne\varnothing$. Immediately before event $j$, every
active row belongs to
\[
 \operatorname{Reach}_P
       \bigl(Z\cup F\cup\{o_1,\ldots,o_{j-1}\}\bigr).
\]
A prerequisite $x\notin C_j$ is not reached from $Z$ or an earlier
output. Reachability from a union is the union of the individual
reachability sets, so some $e\in F$ reaches $x$. This is precisely a
witness in $F\cap D_{j,x}$.

For sufficiency, meeting $A_\ell$ activates every vertex of $Q_\ell$;
hence all of $U$ is available. Induct on the prescribed events. A
prerequisite in $C_j$ is reached from already justified sources or earlier
outputs. Every other prerequisite is reached from a selected edge in its
demand. Thus event $j$ is legal. After the final event every source
component of $P$ contains a singleton seed, a selected edge, or a chosen
output. Every component of a finite acyclic condensation is reachable
from a source component, so ordinary reachability activates every row.
\end{proof}

\subsection{Partitioning the few demands}

The retained demands can overlap and can share a single selected edge.
Requiring a different selected edge for each would be incorrect. The
following elementary lemma isolates the finite combinatorics needed to
handle shared witnesses.

\begin{lemma}[Partition reduction for an independent hitting set]
\label{lem:fpt-partition-reduction}
Let $M$ be a finite matroid on $E$. Let $A_1,\ldots,A_p$ be pairwise
disjoint subsets of $E$, and let $D_1,\ldots,D_q$ be subsets of
$E\setminus\bigcup_\ell A_\ell$. There is an $M$-independent set
meeting every $A_\ell$ and every $D_j$ if and only if, for some set
partition $\pi$ of $\{1,\ldots,q\}$, the family
\[
       A_1,\ldots,A_p,\qquad
       K_B=\bigcap_{j\in B}D_j\quad(B\in\pi)
\]
has an independent system of distinct representatives.
\end{lemma}
\begin{proof}
Given an independent hitting set $F$, choose $e_j\in F\cap D_j$ and
partition the demand indices according to equality of their chosen
witnesses. For every block $B$, its common witness belongs to $K_B$;
witnesses for different blocks are distinct. Also choose a representative
in $F\cap A_\ell$ for each $\ell$. The latter representatives are
pairwise distinct because the $A_\ell$ are disjoint, and none equals a
demand witness because each $D_j$ avoids their union. All representatives
form a subset of the independent set $F$. Conversely, representatives
for the displayed groups meet all original demands and mandatory groups.
The empty demand family uses the unique empty set partition.
\end{proof}

For each partition, solve the independent-transversal problem in the
graphic matroid of $G$. Its groups may overlap. To reduce to the familiar
graphic/partition intersection, give every membership $e\in K$ its own
edge copy $(K,e)$, with the same graph endpoints as $e$. Use the graphic
matroid on these copies and the partition matroid allowing at most one
edge in each required group. Two copies of the same original nonloop
edge are parallel and hence dependent; copies of an original loop remain
loops. A common independent set of size equal to the number of groups
therefore projects to an independent system of distinct original-edge
representatives. There are at most $(q+1)|E|$ copies, not $|E|^q$.

Rado's independent-transversal theorem gives the exact alternative:
either these groups have independent representatives or some subfamily
$\mathcal J$ satisfies
\[
 \rho_G\left(\bigcup_{K\in\mathcal J}K\right)<|\mathcal J|.
 \tag{$\dagger$}\label{eq:fpt-rado-obstruction}
\]
Here the union is a set of original edge labels. Repeated membership does
not increase rank. This is a classical matroid result; the present
reduction supplies the scenario-specific groups to which it applies.

\begin{proof}[Proof of Theorem~\ref{thm:column-fpt}]
Enumerate all scenarios. For each, compute its source groups and at most
$2k$ retained demands. Enumerate their set partitions and solve each
graphic/partition intersection problem. By
Lemma~\ref{lem:fpt-demand-separation},
Lemma~\ref{lem:fpt-partition-reduction} applies; by
Proposition~\ref{prop:fpt-fixed-scenario}, a successful intersection
produces a successful forest. Extend it to a spanning tree. Activation
is monotone in the initially selected rows, so this extension remains
successful; replay peeling directly against the immutable source matrix.

Conversely, any successful tree admits a scenario by
Lemma~\ref{lem:fpt-scenarios-complete}. For that scenario it meets all
required groups and demands. The partition reduction therefore makes
at least one of its matroid-intersection subproblems successful.

There are $3^k k!$ scenarios and at most $B_{2k}$ partitions for each.
All residual constructions and matroid computations are polynomial in
the encoded input size. For $k\ge1$,
\[
      3^k k!B_{2k}
      \le3^k k^k(2k)^{2k}=(12k^3)^k,
\]
which proves the stated bound. Scenarios and partitions can be streamed,
so the producer needs only polynomial working space apart from any
explicitly retained negative certificate.
\end{proof}

\subsection{Independent positive and negative certificates}

A positive certificate consists simply of the selected tree and the
complete peeling schedule. Its checker need not trust, or reproduce,
the scenario search or the matroid-intersection implementation.

For a negative certificate, record a violated inequality
\eqref{eq:fpt-rado-obstruction} for every scenario/partition subproblem.
An independent verifier regenerates the canonical enumeration, derives
each family of requirements, checks that every subproblem is covered,
and verifies the reported rank defect. Graphic rank is the number of
vertices minus the number of connected components of the indicated
spanning subgraph, including isolated vertices. Thus verification uses
only elementary graph operations and the explicit scenario rules; it
does not invoke matroid-intersection optimization. The resulting
certificate is fixed-parameter sized, with the same exponential
parameter factor. It is not asserted to have polynomial size for
unbounded $k$.

All conclusions concern the supplied immutable graph--matrix
certificate class. A negative result does not prove nontriviality of
the knot or exclude other handle systems, slides, subdivisions, or
hierarchy certificates.

\subsection{Optimizing the successful spanning tree}

\begin{theorem}[Minimum-weight successful tree]
\label{thm:fpt-minimum-weight-tree}
Give every eligible edge a binary-encoded signed weight. An exact
minimum-weight successful spanning tree, or proof that none exists,
can be found with the same
$2^{O(k\log(k+2))}\operatorname{poly}(N)$ parameter bound, where $N$
also includes the weight encoding length.
\end{theorem}
\begin{proof}
Fix a scenario and a demand partition, and let $L$ be the number of
required groups. Put $r=|V(G)|-1$. If $L>r$, this subproblem is
infeasible. Otherwise create its required-group copies as above and one
additional dummy copy of each original edge. The first matroid is the
graphic matroid on all copies. The second is a partition matroid with
capacity one for each required group and capacity $r-L$ for the dummy
group. Give every copy the weight of its original edge.

A common independent set of cardinality $r$ saturates every required
group and the dummy capacity. Its projection is a spanning tree
satisfying the group requirements, with exactly the same total weight.
Conversely, a spanning tree having distinct representatives for the
required groups lifts to such a common independent set: mark the
representatives by their group labels and give every remaining tree edge
its dummy label. Weighted matroid intersection therefore gives the exact
minimum tree weight for this subproblem. Minimizing over all scenarios
and partitions gives the global optimum by the completeness proof above.
Signed weights are harmless because the requested cardinality is fixed.
\end{proof}

This weighted theorem is separate from any unweighted reference
implementation. For example, negated row-occurrence counts optimize the
number of occurrences removed initially among successful trees. Such a
weight is a practical surrogate; it does not prove optimality for actual
end-to-end recognition time.

\subsection{Larger exceptional supports and quantitative regimes}

The same proof applies to an exceptional active column with support
$S_i$, size $s_i\ge3$, and nonempty set $U_i\subseteq S_i$ of
unit-coefficient owners. Its output must lie in $U_i$ and its
prerequisites are all other support rows. Put
\[
                         Q=\sum_{i=1}^k(s_i-1).
\]
There are $k!\prod_i|U_i|$ scenarios and at most $Q$ residual demands,
giving the following extension.

\begin{corollary}[Parameterization by nonlinear support]
\label{cor:fpt-nonlinear-support}
Joint forest selection and complete unit peeling with arbitrary
exceptional active columns is decidable in
\[
       k!\left(\prod_i|U_i|\right)B_Q\operatorname{poly}(N)
       =2^{O(Q\log(Q+2))}\operatorname{poly}(N).
\]
The statement remains valid for any matroid with polynomial-time
independence testing in place of the graphic matroid.
\end{corollary}
\begin{proof}
Each exceptional column still has at most one useful activation. Its
scenario event has $s_i-1$ prerequisites, so the preceding arguments
apply verbatim. Since $|U_i|\le s_i$ and $k\le Q$, the displayed
enumeration factor is $2^{O(Q\log(Q+2))}$. Cloning is a parallel
extension of the supplied matroid, and classical matroid intersection
requires only its independence oracle.
\end{proof}

For the triple-unit theorem, $k=O(\log N/\log\log N)$ gives polynomial
time, and $k=O(\log N)$ gives $N^{O(\log\log N)}$. More generally,
\[
                k\log(k+2)=O((\log N)^2)
\]
is sufficient for the quasipolynomial bound $N^{O(\log N)}$; for example,
$k=O((\log N)^2/\log\log N)$ suffices. These implications are
conditional on the corresponding parameter bound. The finite
graph--matrix theorem alone supplies no such uniform bound for all
knot diagrams.

\subsection{Authenticated cap parity preserves rank}

The general algorithm above requires no compatibility between the
occurrence matrix and the graph. The native cap incidence supplies an
additional structure when it has been independently authenticated.
Extend the graphic matroid from $E$ to the whole row set $R$ by making
rows outside $E$ matroid loops. Assume that, for every column of the \emph{full matrix used by activation},
the eligible edges with odd \emph{uncapped} occurrence count form an
even-degree subgraph of the cap graph. This includes every added
finite-boundary or interface column. The distinguished exterior vertex
is included in this condition, and a graph loop contributes degree two.

\begin{theorem}[Rank preservation of positive activation]
\label{thm:fpt-parity-rank}
Under this raw parity hypothesis, every legal unit-target activation
adds a row in the graphic-matroid closure of its prerequisite rows.
For every initial row set $X$,
\[
  \operatorname{cl}_{G}\bigl(\cl_A(X)\bigr)
                 =\operatorname{cl}_{G}(X),
  \qquad
  \rho_G\bigl(\cl_A(X)\bigr)=\rho_G(X).
\]
In particular:
\begin{enumerate}
\item every freely activated row is a graphic-matroid loop;
\item every successful independent seed is already a cap-graph basis;
\item a scenario/partition can have independent representatives only
      when its number $L$ of required groups equals
      $r=\rho_G(E)$.
\end{enumerate}
\end{theorem}
\begin{proof}
Suppose a column activates its unit row $e$. All its other support rows
are already active. If $e$ is a matroid loop, the claim is immediate.
Otherwise $e$ is a nonloop graph edge and, because its coefficient is
one, belongs to the column's odd-occurrence edge set. Every nonloop edge
in an even-degree finite multigraph lies on a cycle: if it were a bridge,
the sum of vertex degrees on one side of that bridge would be odd.
Removing $e$ from such a cycle leaves a path between its endpoints using
other odd-occurrence edges of the same column. These are among the
already active prerequisites. Hence $e$ belongs to their graphic span.
Induction over activations proves equality of the initial and final
graphic spans.

Taking $X=\varnothing$ proves the first consequence. Complete activation
has rank $\rho_G(R)=\rho_G(E)=r$; if $X$ is independent it must therefore have
cardinality $r$, proving the second. For the third, an independent
representative system is a successful independent seed by
Proposition~\ref{prop:fpt-fixed-scenario} and
Lemma~\ref{lem:fpt-partition-reduction}. Its cardinality is exactly the
number of its required groups, so $L=r$.
\end{proof}

The free-closure restriction in
Proposition~\ref{prop:fpt-free-closure} removes only matroid loops under
this authenticated hypothesis, and consequently preserves the raw
even-degree identity. By contrast, replacing an odd multiplicity at
least three by two can change parity. The rank argument must therefore
refer to the original raw counts, although the zero/unit/multiple
predicates used by the decision algorithm can still be capped at two.

\begin{corollary}[One Stirling layer suffices on parity-compatible inputs]
\label{cor:fpt-stirling-layer}
For a fixed scenario with $p$ mandatory source groups and $q$ retained
demands, the parity-compatible algorithm needs only partitions with
\[
                            |\pi|=r-p.
\]
Thus it examines $\left\{\begin{smallmatrix}q\\r-p\end{smallmatrix}\right\}$
partitions, the corresponding Stirling number of the second kind,
instead of all $B_q$ partitions. A scenario with $p>r$ or $p+q<r$
can be rejected immediately. A successful set of representatives is
already a spanning tree when the cap graph is connected.
\end{corollary}
\begin{proof}
There are $L=p+|\pi|$ required groups. The previous theorem forces
$L=r$ in every feasible subproblem. Conversely, excluding the other
partitions loses no solution: a successful tree generates, through the
partition-reduction proof, an independent representative system; rank
preservation forces that system to contain all $r$ of its tree edges.
The usual conventions include one partition of the empty set into zero
blocks and no partition into an invalid number of blocks.
\end{proof}

For a triple-unit column, the earlier source-parity audit's local
classification follows from the same hypothesis \cite{repo-dual-normalization}. Its three row elements form either a triangle of eligible edges, a
parallel pair together with a matroid loop, or three matroid loops, of
graphic ranks two, one, and zero respectively. Indeed, one nonloop edge is incompatible
with even degrees, two nonloop edges must be parallel, and three
nonloop edges must form a triangle. Only the triangle case can obtain
two prerequisites directly from a forest without another activation
route. This classification and rank pruning apply exactly where the
raw chain identity has been verified; they are not inferred merely from
an arbitrary capped matrix.
